\chapter{Matriz de operación de variables}
\label{anexo:d}

La matriz de operación especifica cómo se \textbf{mide} cada ratio del modelo en el diseño del plan.
Complementa la \Cref{tab:variables} (Cap.~2) y la matriz de consistencia (Anexo~C). En PI-11,
\textbf{factores} = condiciones manipuladas; \textbf{ratios} = resultados medidos.

\section{Ratios principales}

\small
\begin{longtable}{@{}p{0.16\textwidth}p{0.10\textwidth}p{0.20\textwidth}p{0.12\textwidth}p{0.22\textwidth}p{0.14\textwidth}@{}}
\caption{Operación de variables de investigación (Anexo D)} \label{tab:anexo-d} \\
\toprule
\textbf{Variable} & \textbf{Fase} & \textbf{Def. operativa} & \textbf{Escala} & \textbf{Instrumento / cálculo} & \textbf{Control} \\
\midrule
\endfirsthead
\multicolumn{6}{l}{\textit{(continúa)}} \\
\midrule
\textbf{Variable} & \textbf{Fase} & \textbf{Def. operativa} & \textbf{Escala} & \textbf{Instrumento / cálculo} & \textbf{Control} \\
\midrule
\endhead
\midrule
\multicolumn{6}{r}{\textit{(continúa en la siguiente página)}} \\
\endfoot
\bottomrule
\endlastfoot
StructuralScore (E) & Fase 3 / OE2 & Correspondencia estructural OpenAPI $\leftrightarrow$ BIAN & [0, 1] continua & Matriz de correspondencias + promedio ponderado de pares alineados & Umbral de matching fijo; mismo insumo normalizado (OE1) \\
SemanticScore (S) & Fase 4 / OE3 & Similitud semántica entre términos OpenAPI y conceptos BIAN & [0, 1] continua & Similitud coseno (embeddings) o medida ontológica acordada & Corpus de términos fijo; modelo de embedding fijo por corrida \\
CoverageScore (C) & Fase 5 / OE4 & Proporción de conceptos BIAN cubiertos por el contrato & [0, 1] continua & Ratio conceptos BIAN mapeados / total de conceptos del SD & Mismos S y E; taxonomía BIAN del SD emparejado \\
AlignmentScore & Fase 5 / OE4 & Indicador global de alineación del contrato & [0, 1] continua & $\alpha \cdot S + \beta \cdot E + \gamma \cdot C$ & $\alpha+\beta+\gamma=1$; umbrales Alta/Media/Baja/Nula fijos \\
Clasificación de alineación & Fase 5 / OE4 & Nivel interpretativo del contrato evaluado & Nominal (4 niveles) & Alta $\geq$0,80; Media 0,60--0,79; Baja 0,40--0,59; Nula $<$0,40 & Mismo AlignmentScore; reglas Tabla tipologías Cap.~1 \\
Calidad de extracción (OE1) & Fase 2 / OE1 & Completitud de la normalización documental & Ratio + conteos & \% elementos extraídos vs.\ artefacto fuente & Versión fija OpenAPI + BIAN; parser fijo \\
\end{longtable}

\normalsize

\section{Factores por objetivo específico}

\begin{table}[H]
  \centering
  \caption{Operacionalización de factores}
  \label{tab:anexo-d-vi}
  \small
  \begin{tabular}{@{}p{0.08\textwidth}p{0.24\textwidth}p{0.58\textwidth}@{}}
    \toprule
    \textbf{OE} & \textbf{Factor} & \textbf{Estados / niveles considerados en el diseño} \\
    \midrule
    OE1 &
    Esquema de representación intermedia &
    Tabular (recurso--operación--esquema); grafo (nodos/aristas); RDF/triples (opcional) \\
    OE2 &
    Estrategia de \textit{schema matching} &
    Nombre exacto; similitud de ruta; profundidad de esquema; operación HTTP \\
    OE3 &
    Técnica de similitud semántica &
    Embeddings + coseno; matching ontología; híbrido estructural+semántico \\
    OE4 &
    Pesos $\alpha$, $\beta$, $\gamma$ &
    Pesos iguales ($1/3$); ponderación experta; análisis de sensibilidad \\
    \bottomrule
  \end{tabular}
\end{table}

\section{Notas de medición}

\begin{itemize}
  \item Todas las ratios de score se normalizan al intervalo $[0, 1]$ para permitir agregación
        ponderada.
  \item La clasificación nominal (Alta/Media/Baja/Nula) se obtiene por umbrales fijos definidos en
        el Capítulo~1 §1.4 y reutilizados en OE4.
  \item El diseño documenta el procedimiento; la ejecución experimental a escala queda fuera de
        alcance del plan (Guía~N°01, §1.6).
\end{itemize}
